\documentclass[12pt]{article}
\usepackage[utf8]{inputenc}
\usepackage[T1]{fontenc}
\usepackage{hyperref}
\title{Enawen\'e-Naw\'e Language Resources for K-12 Curriculum Development}
\author{Compiled by Claude Sonnet 5 (claude-sonnet-5) for Andr\'as Kornai}
\date{}
\begin{document}
\maketitle
\section{Introduction}
Enawen\'e-Naw\'e (ISO 639-3 unk; Glottocode enaw1238) is an Arawakan (Maipurean) language of Mato Grosso, Brazil. Sources report strong intergenerational transmission, and most of the people are still monolingual. The search found no native-language textbooks or primers.
\section{Primary Educational Resources}
None found. A 2019 legislative request asked for work to resume on the Enawene Nawe school in the village of Halataikwa, Ju\'ina (\url{https://www.al.mt.gov.br/storage/webdisco/cp/20190716100239155600.pdf}). I did not find the school's current status.
\section{Linguistic and Reference Materials}
A review of Katia Zorth\^ea's work on the acquisition of alphabetic writing by the Enawen\'e-Naw\'e, which says it happened outside the school (\url{https://revistapedagogia.ufmt.br/wp-content/uploads/2025/08/resenha_Ana-Paula-06-11.pdf}). A UEMS master's dissertation on their schooling at CEJA Alternativo, Ju\'ina (\url{https://biblioteca.uems.br/publicos/download/96444}). Endangered Languages Project entry: \url{https://endangeredlanguages.com/lang/2054}.
\section{Religious and Community Resources}
Povos Ind\'igenas no Brasil profile: \url{https://indios.org.br/en/povo/Enawen\%C3\%AA-naw\%C3\%AA}.
\section{Digital and Audio Resources}
ELDP-archived ``Language Documentation of Enawene-Nawe (Arawak)'': \url{https://eldp.access.preservica.com/dk0572/}.
\section{Curriculum Development Context}
The language is not known to be taught in schools, though the school at Halataikwa has been under discussion. The ELDP recordings and the academic studies above are the most relevant starting points. Any curriculum work needs community consent.
\end{document}